\cleardoublepage
\chapternonum{致谢}

十分感谢刘勇老师与李智彬老师，在他们的教导与帮助下，我理解了撰写论文的初衷与科研的意义，我的各种能力得到了较大的提升。十分感谢朱承睿师兄与张震师兄，在与他们的交流过程中，我的工程实践能力得到了较大的提升。十分感谢朱德烨、翟光耀、李宝润、陈炳燃、李庆鹏、吴艳明、袁昕、游梁成、段兴宇、康宇鑫等师兄弟们，我们一起探讨机器人技术的发展与学习方法的理论，并肩进步。十分感谢彭林鹏、梅剑标、刘维维、杜思怡、丁媛缘、傅晨灿、陈俊豪、蔡荣耀、郭闯、徐李明、胡天扬、王龙祺、韦佳腾、芦正明、李胜博、刘馨阳等 APRIL 实验室的同学们，与你们共同度过了这段愉快的研究生时光。最后，十分感谢我的父母，我的妹妹马林馨，我的女朋友胡懿芝与我的同学王安家、张庆普、樊泽华等人，你们始终提供鼓励、理解与支持，让我不断前行。

哈工大教会了我“规格严格”，浙大教会了我“求是创新”，这些都是令我受益匪浅的人生道理。当前机器人技术虽发展迅速，但距离广泛的应用仍有较长的路要走。在这一过程中，我感觉World Model将是一个至关重要的方法，当然，这仅仅是个人观点。在游戏、视频与音频领域，生成式世界模型依然占据主流；但在机器人控制领域，Latent Manifold World Model可能成为关键。未来的范式可能是由生成式模型提供目标状态轨迹，再由Latent Manifold World Model提供规划器以恢复出最终策略。无论是从生物学启示还是工程实践的角度来看，世界模型都更像是处理多模态、多任务及多场景场景下机器人控制问题的正确途径。

当然上面都是个人观点，我的这篇论文没有相关论述。这篇论文怎么说呢，提供的思想也比较简单，可能也没多少人看。不过我可能也努力了。这三年，想一想，还是挺充实的，应该也算是我从生下来进步最大的三年。我的毕业是一个阶段的结束，也是一个阶段的开始。我猜2026将会是Agent与世界模型爆发的一年，AI与机器人领域将会有很大的改变。最后十分感谢April这个包容的大家庭，让我有机会与这么多优秀的同学们一起学习与成长。同时祝看到这篇论文的人一切顺利。
